% Clase 5 --- Prisma triangular en campo uniforme: flujo neto cero (§5).
% La cara inclinada tiene mayor área que la frontal, pero su proyección
% perpendicular al campo es EXACTAMENTE A1: A3 cos(theta) = A1. Por eso el
% flujo saliente compensa al entrante.
\begin{center}
\begin{tikzpicture}[scale=1]

  % campo uniforme
  \foreach \y in {-1.2,-0.6,0,0.6,1.2} {
    \draw[figvecres] (-3.6,\y) -- (-2.3,\y);
  }
  \node[figlblaux, above=2pt] at (-2.95,1.25) {$\vec v$ uniforme};

  % prisma visto de canto (triángulo) + perspectiva
  \draw[figline, fill=figblue!6] (-1.4,-1.4) -- (-1.4,1.4) -- (1.7,-1.4) -- cycle;
  \draw[figaux] (-1.4,1.4) -- (-0.75,1.85) -- (2.35,-0.95) -- (1.7,-1.4);
  \draw[figaux] (-1.4,-1.4) -- (-0.75,-0.95) -- (2.35,-0.95);
  \draw[figaux] (-0.75,1.85) -- (-0.75,-0.95);

  % cara frontal A1 y su normal
  \node[figlbl, figblue, above left=1pt] at (-1.42,0.55) {$A_1$};
  \draw[figvec] (-1.4,0.3) -- (-2.15,0.3);
  \node[figlblaux, left=2pt] at (-2.15,0.3) {$\hat n_1$};

  % cara inclinada A3 y su normal
  \node[figlbl, figblue, right=3pt] at (0.55,-0.45) {$A_3$};
  \draw[figvec] (0.15,0.0) -- (0.83,0.75);
  \node[figlblaux, right=2pt] at (0.83,0.75) {$\hat n_3$};
  % ángulo entre n3 y el campo
  \draw[figvecaux] (0.15,0) -- (0.95,0);
  \draw[figaux] (0.72,0) arc (0:48:0.57);
  \node[figlblaux] at (26:0.95) {$\theta$};

  % cara superior A2
  \node[figlbl, figblue, above=1pt] at (-1.1,1.55) {$A_2$};
  \draw[figvec] (-1.1,1.62) -- (-1.1,2.3);
  \node[figlblaux, above=1pt] at (-1.1,2.3) {$\hat n_2$};

  % balance de flujos
  \node[figlbl, align=left, anchor=west] at (3.2,0.75)
    {$\Phi_1 = -v A_1$\\[3pt]
     $\Phi_3 = +v A_3\cos\theta = +v A_1$\\[3pt]
     $\Phi_2 = \Phi_4 = \Phi_5 = 0$};
  \node[figlbl, figred, anchor=west] at (3.2,-0.85)
    {$\Phi_{\text{total}} = 0$};

\end{tikzpicture}\\[0.5em]
{\small\itshape Las caras paralelas al campo no aportan nada ($\vec v \perp
\hat n$). La cara inclinada es \textbf{más grande} que la frontal, pero el
$\cos\theta$ la corrige exactamente: $A_3\cos\theta = A_1$. Entran tantas líneas
como salen porque \textbf{no hay carga adentro} ---ni fuentes ni sumideros---, y
ése es justamente el germen de la ley de Gauss.}
\end{center}
